% SPDX-License-Identifier: LPPL-1.3c
% Copyright (C) 2026 Christophe Jorssen
% Author and Current Maintainer: Christophe Jorssen <christophe.jorssen@gmail.com>
% LPPL maintenance status: maintained
% This file is part of luafemm. See LICENSE and LICENSES.md for its terms.
\documentclass[varwidth=150mm,border=3mm]{standalone}
\usepackage{tikz}
\usetikzlibrary{femm}
\begin{document}
\tikzset{femm/materials/eccentric iron/.style={mu r=100}}
\begin{tikzpicture}[
  femm/problem={xmin=-8,xmax=8,ymin=-8,ymax=8,mesh size=1,depth=10},
  femm/sliding band={name=eccentric-ring,inner radius=3,
    outer radius=4,segments=64,layers=2},
  scale=2.5,baseline=(current bounding box.center)]
  \draw[femm/region={material=eccentric iron}]
    (0,0) ellipse[x radius=2.6,y radius=1.2];
  \draw[femm/region={material=copper,current density=1e7}]
    (-6,-1) rectangle (-5,1);
  \draw[femm/region={material=copper,current density=-1e7}]
    (5,-1) rectangle (6,1);
\end{tikzpicture}
\begin{tikzpicture}[
  femm/problem={sliding band=eccentric-ring,rotor angle=25,
    rotor offset x=.45,rotor offset y=.2},
  scale=2.5,baseline=(current bounding box.center)]
  \pic {femm geometry};
  \pic[femm/lines=13] {femm field};
  \draw[teal,dashed] (.45,.2) circle[radius=3];
  \draw[teal,dashed] (0,0) circle[radius=4];
  \femmmeasure[name=eccentric-force,quantity=force,method=band,
    origin x=.45,origin y=.2]
\end{tikzpicture}
\end{document}
